\begin{lemma}{Eigenschaften Gordon'scher Paare}
             {EigenschaftenGordonscherPaare}
Seien $c \in \R$ und $(X, Y)$ das Gordon'sche Paar zu $(\gamma, \delta,g,c)$.\\
Der Raum $\R^m \times \R^N$ sei mit der Produktnorm ausgestattet.\\
Dann gilt:
\begin{enumerate}[(1)]
	\item \label{EigenschaftenGordonscherPaare1}
	$\forall\, t \in \R^m\, \forall\, v \in \R^N:
	\E(X(t,v)) = 0 = \E(Y(t,v))$.
	\item \label{EigenschaftenGordonscherPaare2}
	$\forall\, t \in \R^m\, \forall\, v \in \R^N: \norm{X(t,v)}_2^2 
	= c^2\norm t _{2,m}^2 + \norm v _{2, N}^2$.
	\item \label{EigenschaftenGordonscherPaare3}
	$\forall\, t \in \R^m\, \forall\, v\in \R^N:
	\norm{Y(t, v)}_2^2 = \norm t_{2,m}^2\norm v_{2, N}^2$.
	\item \label{EigenschaftenGordonscherPaare4}
	$X$ ist linear und stetig.
	\item\label{EigenschaftenGordonscherPaare5}
	$Y$ ist bilinear und stetig.
\end{enumerate}
\end{lemma}
\begin{proof}
(1) Seien $t \in \R^m$ und $v \in \R^N$.\\
Falls $ct\neq 0$ oder $v \neq 0$, ist $X(t,v)$ nach
\refclaim{IntegrierbarkeitUndVerteilungVonLinearkombinationen}
         {IntegrierbarkeitUndVerteilungVonLinearkombinationen4}
$N(0, c^2\norm t_{2, m}^2 + \norm v_{2,N}^2)-$verteilt, hat also Erwartungswert
0 und Varianz $c^2\norm t_{2, m}^2 + \norm v_{2,N}^2$ 
(siehe 8.8, 9.2 in~\cite{Irle}).\\
Ist $ct = 0$ und $v = 0$, so ist $X(t, v) = 0$ und hat daher trivialerweise den
Erwartungswert 0 und auch Varianz 
$0 = c^2\norm t_{2, m}^2 + \norm v_{2,N}^2$.\\
Somit ist $\E(X(t, v)) = 0$ und daher
$\norm{X(t,v)}_2^2 = c^2\norm t_{2, m}^2 + \norm v_{2,N}^2$ (vgl. Bemerkung vor~\ref{ElementareEigenschaftenQuadratintegrierbarerAbbildungen}).\\
Falls $t \neq 0$ und $v \neq 0$, so ist $Y(t,v)$ nach~\refclaim{IntegrierbarkeitUndVerteilungVonLinearkombinationen}
         {IntegrierbarkeitUndVerteilungVonLinearkombinationen5}
$N(0, \norm t _{2,m}^2 \norm v_{2, N}^2)-$verteilt, hat also Erwartungswert 0
und Varianz $\norm t_{2,m}^2\norm v_{2,N}^2$.\\
Falls $t = 0$ oder $v = 0$, so ist $Y(t, v) = 0$ und hat wiederum trivial
Erwartungswert 0 und Varianz $0=\norm t _{2,m}^2 \norm v_{2, N}^2$.\\
Also ist $E(Y(t, v))= 0$ und 
somit $\norm{Y(t, v)}_2^2 = \norm t _{2,m}^2 \norm v_{2, N}^2$ (mit dem
gleichen Argument, wie bei $X(t, v)$).\\
Damit sind (1), (2) und (3) gezeigt.\\
(4) Nach~\refclaim{EigenschaftenDesLinearkombinators}
{EigenschaftenDesLinearkombinators1}
sind \[c\bprod\cdot \delta_m:\R^m \rightarrow L_2(\Omega,\mfa, \Prim, \R)
\quad \text{und}\quad 
\bprod \cdot \gamma_N:\R^N \rightarrow L_2(\Omega,\mfa, \Prim,\R)\] linear.
Daher ist deren Summe linear im Tupel der beiden Komponenten.\\
Als lineare Abbildung mit endlichdimensionalem Definitionsraum ist $X$
stetig.\\
(5) Seien $t,s\in \R^m$, $v, w \in \R^N$, $a \in \R$ und $\omega \in \Omega$.\\
Dann gilt wegen der Linearität von $g(\omega)$:
\begin{align*}
Y(t+as, v)(\omega)& = \sprod{g(\omega)(t+as)}v_N 
= \sprod{g(\omega)t + ag(\omega)s}v_N\\
&= \sprod{g(\omega)t}v_N + a\sprod{g(\omega)s}v_N =
Y(t, v)(\omega)+ a Y(s, v)(\omega)\\
&= (Y(t, v) + a Y (s, v))(\omega)\text.
\end{align*}
Außerdem ist
\begin{align*}
Y(t, v+aw)(\omega) &= \sprod{g(\omega)t}{v+aw}_N
= \sprod{g(\omega)t}v_N+a\sprod{g(\omega)t}w_N\\
&= Y(t, v)(\omega) + a Y(t, w)(\omega)
= (Y(t, v) + a Y(t, w))(\omega)\text.
\end{align*}
Somit ist $Y$ bilinear.\\
Wegen (3) ist $\norm{Y(t, v)}_2 = \norm t_{2,m}\norm v_{2,N}$ und da $Y$
bilinear ist, ist es nach der bekannten Charakterisierung der Stetigkeit
bilinearer Abbildungen stetig. Auch hier könnte man alternativ die endliche
Dimension des Definitionsraums nutzen, um die Stetigkeit von $Y$ einzusehen:
$\R^m, \R^N$ und $L_2(\Omega, \mfa, \Prim, \R)$ sind vollständig, $Y$ ist
komponentenweise als lineare Abbildung mit endlichdimensionalem
Definitionsraum stetig, also ist nach dem Satz von Banach--Steinhaus auch
$Y$ stetig.
\end{proof}